\section{Conclusions}

This project has demonstrated that a functional stack-based virtual machine 
complete with a custom instruction set, an assembler, and an interactive REPL 
can be implemented in roughly 600 lines of C with no dependencies beyond the 
standard library. The result is a system small enough to be read and understood 
in its entirety, yet expressive enough to evaluate arithmetic expressions, 
implement conditional logic, execute recursive function calls, and maintain 
persistent state across loop iterations using sixteen general-purpose registers.

Several design decisions proved their value through the implementation process. 
Centralising \texttt{push} and \texttt{pop} in \texttt{vm.c} ensured that 
stack-safety checks exist in exactly one place each, eliminating an entire 
class of potential inconsistency. Separating the operand and call stacks 
prevented return addresses from being silently corrupted by arithmetic 
operations — a subtle correctness property that real processors enforce in 
hardware. Zero-initialising the \texttt{VM} struct with \texttt{= \{0\}} 
eliminated a category of non-deterministic behaviour at negligible cost. And 
confining each opcode category to its own translation unit kept the codebase 
navigable throughout development.

The implementation notes in Section~5 are equally instructive for what they 
reveal about the cost of incomplete error handling. The latent bugs in 
\texttt{op\_div} and \texttt{op\_store}, both caused by missing \texttt{return} 
statements after setting \texttt{running = 0} – illustrate a recurring pitfall 
in C: the language provides no mechanism to enforce that an error branch 
terminates. In a larger system these defects could manifest as memory corruption 
or undefined behaviour that is difficult to reproduce. The correct pattern, 
demonstrated in every other handler, is to use early returns and 
\texttt{else} guards so that error paths and computation paths are structurally 
mutually exclusive.

The comparison with the JVM, CPython, and WebAssembly in Section~6 places 
this project in a wider context. All four systems share the same stack-based 
execution model, which confirms that the stack machine is not merely a 
pedagogical simplification — it is an architectural choice that production 
language designers have independently adopted because its specification and 
implementation costs are genuinely low. What distinguishes the production 
systems is the machinery layered on top: type systems, bytecode verifiers, 
garbage collectors, just-in-time compilers, and formally specified memory 
models. This VM occupies the simplest possible position on every one of those 
axes, and that simplicity is the point. Every layer that was deliberately 
omitted is a layer whose absence makes the remaining design easier to reason 
about.

Future work could extend the VM in several directions. Adding a linear 
byte-addressable memory region analogous to WebAssembly's linear memory 
would allow the VM to represent strings and arrays, removing the most 
significant expressive limitation of the current design. Introducing a type tag 
in the high bits of each stack word would enable basic runtime type checking 
without changing the integer representation of values. Bounds-checking jump 
targets before assigning to \texttt{ip} would close the only remaining source 
of undefined behaviour. Each of these extensions can be made incrementally, 
touching at most the files identified in the separation-of-concerns analysis, 
which validates the modularity of the original architecture.

The virtual machine would also benefit from save/load using .txt files in order to save
previously-written programs; as well as a separate interface using C++ library like SFML
to provide cleaner visual identity and editing capabilities to the software.

Ultimately, this project achieves the goal stated in its motivation: to make 
explicit every design decision that a language runtime ordinarily hides. The 
fetch–decode–execute loop, the implicit operand convention of a stack machine, 
the mechanics of a call stack, and the encoding of a flat integer bytecode 
stream are no longer abstractions. They are concrete, readable, and — 
intentional bugs aside — correct.